% Building 2, Part II. The edge.
% One file for the whole part. Each chapter begins with a banner and \section*.

\part*{Part II. The edge}
\phantomsection\addcontentsline{toc}{part}{Part II. The edge}

Where the method stops proving and starts sorting. Each chapter states one line that is not yet among the rules of Building 0, Chapter 0 (they enter there only when agreed), runs it on the map, and says honestly what is proved, what is checked by machine, and what is reported. Nothing here is a new theorem about the primes.

%% ==================================================================
%%  Chapter 3  (loops)
%% ==================================================================
\section*{Chapter 3. Loops}
\phantomsection\addcontentsline{toc}{section}{\texorpdfstring{Chapter 3. Loops \quad a later system blurred by an earlier one $\to$ what comes back}{Chapter 3. Loops}}
\begin{center}
{\normalsize a later system blurred by an earlier one $\to$ what comes back \quad$\cdot$\quad from: Building 0, Chapters 8, 16, 18, 19, 22, II.6; Building 1, Chapters II.5, II.12 \quad$\cdot$\quad blur}\\[4pt]
{\small Build something late, let something early act on it, and blur by the action. Sometimes the early system comes back exactly. Sometimes nearness dies and the bags cannot be told apart: then keep the moves. One prime alone is tame; two primes together are already wild.}
\end{center}

\subsection*{The line of this chapter}

\begin{lines}
\item[\lab{L}] (Loop; draft rule F8.) A chapter $K$ builds a system $X$ and blurs it by the action of an earlier system $S$ (an orbit blur, Building 1, Chapter II.5). A \emph{loop} is a theorem of $K$ that $S$ comes back: a question on the bags whose answers match $S$, respecting the operations named. Verdicts: \emph{closed} ($S$ comes back exactly); \emph{changed} (larger or smaller); \emph{wild} (the bags are not kept apart by any rung of nearness: keep the moves, do not forget them). No chapter before $K$ may use the loop (F7).
\end{lines}

\subsection*{Propositions: the loops on the map}

\begin{lines}
\item[\lab{P1}] (Closed: $\Z$ on $\R$.) Blur $\R$ by $\Z$ acting by $+$ (Building 0, Chapter 8). The turns of loops in $\R/\Z$ are whole numbers, and every whole number occurs (Building 0, Chapter 16, T6): $\Z$ comes back.
\item[\lab{P2}] (Tame: one prime.) Let the powers of $p$ act on $(0, \infty)$ by $\times$, and read through $\log$ (Building 0, Chapter 22): the bags are a circle of length $\log p$, a picture of Building 0, Chapter 8. Likewise the units of $\Z_p$ acting on $\Z/p^k$ by $\times$ have exactly $k + 1$ bags, one per valuation (Building 0, Chapter 19).
\item[\lab{P3}] (Wild: two primes.) Let $2$ and $3$ act on $(0, \infty)$ by $\times$. Read through $\log$, the orbit of $1$ is $\{a \log 2 + b \log 3\}$, and $\log 3 / \log 2$ is not a fraction ($2^a = 3^b$ has no answer with $a, b \ge 1$, by unique factorisation). So the orbit is dense (Building 0, Chapter 8, T7, irrational turns): every rung of nearness glues every bag to every other. Nearness dies; the loop is wild. Machine: the smallest gap among $a \log 2 + b \log 3 \le L$ falls as $L$ grows ($L = 10, 20, 40, 80$).
\item[\lab{P4}] (Closed, with the moves kept: $\Z/p$.) On $\Z/p$, $+$ gives the waves $\psi(x) = e^{2\pi i x/p}$ and $\times$ gives the characters $\chi$ of the clock $(\Z/p)^\times$ (Building 1, Chapter II.12, T5). Mix them: $\tau(\chi) := \sum_x \chi(x) \psi(x)$. Then for every $\chi \ne 1$ \[ |\tau(\chi)|^2 = p . \] (Write $|\tau|^2 = \sum_{x, y} \chi(x)\overline{\chi(y)} \psi(x - y)$; put $x = ty$ and sum over $y$ first: the sum of a wave over a clock is $p - 1$ or $-1$; then $\sum_t \chi(t) = 0$.) A size comes free from $+$ (Parseval) and $\chi$ respecting $\times$.
\item[\lab{P5}] (Lands on a known line: $\det$ acting on ``divides''.) The $n \times n$ matrix $R$ with $R_{ij} = 1$ when $j = 1$ or $i \mid j$ has $\det R = M(n) = \sum_{k \le n} \mu(k)$ (classical, Redheffer; checked $n \le 40$). The line $|M(n)| \le C_\varepsilon n^{1/2 + \varepsilon}$ is equivalent to the Riemann hypothesis (classical, reported).
\end{lines}

\subsection*{Theorems}

\begin{lines}
\item[\lab{T1}] (What comes back, sorted.) On the map: loops by $+$ give a mirror (P1; on every closing-up of $\Q$ at once, $\Q$ comes back as its own waves: reported, Tate); loops by $\times$ give the primes (P2; reported, class field theory for $\Q$); the loop mixing them, $\Q^\times$ acting on all closings-up with $+$, is wild (P3 is its smallest piece: reported, Connes).
\item[\lab{T2}] (Wildness is not local.) Every single prime gives a tame loop (P2). Two primes already give a wild one (P3). So no object built one prime at a time can carry a size on the mixed loop.
\end{lines}

\subsection*{Remarks}

\begin{lines}
\item[\lab{R1}] (What the method did.) The loop line sorts what is on the map: closed, tame, wild. It proves P1--P4 and T2; T1's three verdicts are known mathematics, placed.
\item[\lab{R2}] (Finite floors, moves kept.) On curves over $\Z/p$ every count's error is a sum of mixed sums like P4, each of size $\sqrt p$ (Jacobi sums; checked for $x^3 + y^3 = 1$, $20$ primes). On $\Z$ the same mixed sums appear in the mirror of a series, not in its zeros.
\item[\lab{R3}] (Old names.) P4: Gauss sums; T1: Tate's thesis (self-duality of the adeles), class field theory, the adele class space; P5: the Redheffer matrix, the Mertens function.
\item[\lab{R4}] (Machine check, \texttt{checks/b2-p2-check.py}.) P4 for every character mod every prime $p < 80$; P2, P3 gaps; P2's $k + 1$ bags for $2^5$, $3^4$, $5^3$; P5 for $n \le 40$.
\end{lines}

\subsection*{Exercises}

\begin{lines}
\item[\lab{E1}] Compute $\tau(\chi)$ for the sign character mod $5$ ($\chi(x) = \pm 1$ by squares).
\item[\lab{E2}] Why is the loop of $2$ and $4$ acting by $\times$ tame?
\end{lines}

\newpage

%% ==================================================================
%%  Chapter 4  (amplify)
%% ==================================================================
\section*{Chapter 4. Amplify}
\phantomsection\addcontentsline{toc}{section}{\texorpdfstring{Chapter 4. Amplify \quad a bound on every rung $\to$ a sharp bound}{Chapter 4. Amplify}}
\begin{center}
{\normalsize a bound on every rung $\to$ a sharp bound \quad$\cdot$\quad from: Building 0, Chapters 5, 17, 20, 28, 34; Building 1, Chapters II.12, III.3 \quad$\cdot$\quad ladder, blur}\\[4pt]
{\small Blur, moves and loops ask ``same or different?''. Amplify asks ``how much?''. Climb the powers of one move, bound every rung with any constant, and blur sequences by growth: the constant dies and the bound becomes sharp.}
\end{center}

\subsection*{The line of this chapter}

\begin{lines}
\item[\lab{L}] (Amplify; draft rule Z12.) Given a move, its power ladder, and rung counts $S_m = \sum_i c_i \alpha_i^m$ with finitely many distinct $\alpha_i$ and every $c_i \ne 0$. Blur sequences by growth: $a \preceq b$ when $\limsup |a_m|^{1/m} \le \limsup |b_m|^{1/m}$ (one-way; Building 0, Chapter 17). \\
\emph{Survives:} ``the largest $|\alpha_i|$''. So if $|S_m| \le C (m + 1)^d \rho^m$ for every $m$, with any $C$ and $d$, then every $|\alpha_i| \le \rho$. \\
\emph{Dies:} the constant $C$, and $(m + 1)^d$ ($C\rho^m$ and $\rho^m$ are glued).
\item[\lab{T1}] (Proof.) The bound makes $\sum_m S_m t^m$ converge for $|t| < 1/\rho$ (Building 0, Chapter 20, comparison). That sum is $\sum_i c_i \alpha_i t / (1 - \alpha_i t)$, whose poles are exactly the $1/\alpha_i$ (distinct, no $c_i = 0$); a series that is a fraction of polynomials converges exactly out to its nearest pole (Building 0, Chapter 34, T1 and Chapter 33, P2). So no $1/\alpha_i$ is inside $|t| < 1/\rho$.
\end{lines}

\subsection*{Propositions: three examples}

\begin{lines}
\item[\lab{P1}] (Cauchy--Schwarz with constant $1$.) Suppose only that $|\langle x, y\rangle| \le C |x| |y|$ for all lists of all lengths, one $C$ for all. Apply it to $a^{\otimes k}$, $b^{\otimes k}$ (Building 1, Chapter II.8, T3): $\langle a^{\otimes k}, b^{\otimes k}\rangle = \langle a, b\rangle^k$ and $|a^{\otimes k}| = |a|^k$, so $|\langle a, b\rangle| \le C^{1/k} |a| |b|$ for every $k$, and $C^{1/k}$ comes near $1$. \ Dies when the constant depends on the length: $n^{k}$ on rung $k$ gives back $n$.
\item[\lab{P2}] (The trace finds the largest eigenvalue.) For a square matrix $A$ with eigenvalues $\lambda_i$, $\operatorname{tr} A^m = \sum \lambda_i^m$ (each with its multiplicity, never $0$), and $\lim |\operatorname{tr} A^m|^{1/m}$ along a suitable part is $\max |\lambda_i|$ (T1; one entry of $A^m$ can hide an eigenvalue, the trace cannot).
\item[\lab{P3}] (The road to Building 1, Chapter III.3, T3.) Move: $x \mapsto x^p$ on a curve over $\mathbb{F}_{p^m}$ (Building 1, Chapter II.12). Rung $m$: $S_m = p^m + 1 - N_{p^m}$. If $S_m = \alpha_1^m + \dots + \alpha_{2g}^m$ with $\alpha_i \alpha_{i+g} = p$ (reported: rationality and the mirror) and $|S_m| \le C p^{m/2}$ for every $m$ (reported: a rough bound by an auxiliary polynomial), then $|\alpha_i| \le \sqrt p$ by L, and $= \sqrt p$ by the mirror; with $g = 1$, $|N_p - p - 1| \le 2\sqrt p$. Two inputs reported; Amplify itself is proved.
\end{lines}

\subsection*{Theorems}

\begin{lines}
\item[\lab{T2}] (What Amplify never does.) It removes losses that do not grow with the rung. A loss of a fixed factor per rung ($C^m$) is not removed: it changes $\rho$.
\item[\lab{T3}] (On $\Z$, one slot is empty.) For the zeros of the series $\zeta$ there are power ladders with a Close step (reported: Li's ladder, $s \mapsto 1 - 1/s$; the cut ladder $x \mapsto x^m$ in $\psi(x) - x$; convolution powers of a test map). What no one has is the weak bound: a bound on every rung losing only a constant. Every known bound loses a fixed factor per rung (T2).
\end{lines}

\subsection*{Remarks}

\begin{lines}
\item[\lab{R1}] (What the method did.) Amplify is built from parts already in the books: a move (Building 0, Chapter 0, W2), a ladder (Chapter 5), the growth blur (Chapter 17), and a Close step. Its proof is T1. The three proofs of the curve line (Stepanov--Bombieri, Weil, Deligne) are one chain read this way. Known mathematics (amplification, the tensor power trick); its place in the grammar is the book's.
\item[\lab{R2}] (Finite floors, checked.) The genus-$2$ curve $y^2 = x^5 + 2x + 1$ over $\mathbb{F}_{3^m}$: $N = 7, 15, 19, 83$ for $m = 1$ to $4$; the four $\alpha_i$ from $N_1$, $N_2$ all have size $\sqrt3$ and predict $N_3$, $N_4$ exactly. Graphs: a ladder whose weak bound fails is not saved; the bound is the whole content.
\item[\lab{R3}] (Old names.) L: amplification; P1: Cauchy--Schwarz by tensor powers; P2: Gelfand's formula (trace form); P3: Hasse's theorem by Weil's route; T3: Li's criterion (Li; Bombieri--Lagarias).
\item[\lab{R4}] (Machine check, \texttt{checks/b2-p2-check.py}.) P1's identities for $k \le 4$; P2 on a random $5 \times 5$ matrix; R2's counts in $\mathbb{F}_{3^m}$, $m \le 4$, and the predictions.
\end{lines}

\subsection*{Exercises}

\begin{lines}
\item[\lab{E1}] Use P1's route to show $|\sum a_i b_i c_i| \le (\sum a_i^3)^{1/3} (\sum b_i^3)^{1/3} (\sum c_i^3)^{1/3}$ for $a_i, b_i, c_i \ge 0$, given it with some constant.
\item[\lab{E2}] Why does P2 fail for one entry of $A^m$? Find a $2 \times 2$ example.
\end{lines}

\newpage

%% ==================================================================
%%  Chapter 5  (atoms)
%% ==================================================================
\section*{Chapter 5. Atoms}
\phantomsection\addcontentsline{toc}{section}{\texorpdfstring{Chapter 5. Atoms \quad things built by a glue $\to$ their bags of atoms}{Chapter 5. Atoms}}
\begin{center}
{\normalsize things built by a glue $\to$ their bags of atoms \quad$\cdot$\quad from: Building 0, Chapter 18; Building 1, Chapters II.3, II.7, II.14 \quad$\cdot$\quad blur, ladder}\\[4pt]
{\small Forget how a thing was glued from atoms and keep only the bag of atoms. Sometimes the bag tells everything, sometimes only its size, sometimes nothing. And a floor with two glues is where the hard lines live.}
\end{center}

\subsection*{The blur}

On a floor with one glue (an operation) whose members are built from atoms (members that cannot be glued from smaller ones). \ B1: forget the order and the brackets of the gluing; keep the bag of atoms.

\subsection*{Propositions: a ladder of floors}

\begin{lines}
\item[\lab{P1}] (Rung 1, free: one bag builds one thing.) $\Z$ under $\times$ (Building 0, Chapter 18, T3); $K[x]$ (Building 1, Chapter II.7, T2); $\Z[i]$.
\item[\lab{P2}] (Rung 2, fair: the bag of rungs survives, the glue does not.) Finite move systems, glued by ladders of blurs: the list of simple rungs is kept (Building 1, Chapter II.14, T3), but $\Z/4$ and $\Z/2 \times \Z/2$ have the same bag. $\Z[\sqrt{-5}]$: factorisations differ (Building 1, Chapter II.7, P3) but every factorisation of a member has the same number of atoms (classical, Carlitz; checked to size $600$).
\item[\lab{P3}] (Rung 3, unfair: not even the number survives.) $\Z[\sqrt{-14}]$: $3$ and $5 \pm 2\sqrt{-14}$ are atoms (no member has size $3$ or $27$), and \[ 81 = 3 \cdot 3 \cdot 3 \cdot 3 = (5 + 2\sqrt{-14})(5 - 2\sqrt{-14}) , \] four atoms against two.
\item[\lab{P4}] (Rung 4, atomless.) $\Z$ under $+$: $\Z \supset 2\Z \supset 4\Z \supset \cdots$ never ends; $\Q$, $\R$. Use generators or nearness instead.
\item[\lab{P5}] (Rung 5, uncomputable.) Move systems given by words and gluings: whether two words are glued cannot be decided in general (Building 1, Chapter 2, T4, reported).
\end{lines}

\subsection*{Floors with two glues}

\begin{lines}
\item[\lab{P6}] (The table.) \\
$\Z$: free under $\times$, atomless under $+$; $+$ does not respect the primes; the line about primes (Chapter 6) open. \\
Polynomials over $\Z/p$: free under $\times$; $+$ has atoms ($p \cdot 1 = 0$); irreducibles counted exactly (Building 1, Chapter II.12, T8). \\
Curves over $\Z/p$: fair under $\times$; $+$ has atoms; a move ($x \mapsto x^p$) respects both glues; the line proved (Chapter 4, P3). \\
$(0, \infty)$: atomless under both, and $\log$ turns one into the other: nothing to prove. \\
Atoms with sizes but no exact $+$ (Beurling's primes): the line can fail (reported).
\item[\lab{P7}] (Observation, not a theorem.) The line is proved where $+$ has atoms and one move respects both glues; it is open where $\times$ is free and $+$ is atomless; with no exact $+$ it can fail.
\end{lines}

\subsection*{Remarks}

\begin{lines}
\item[\lab{R1}] (What the method did.) One blur (forget the glue) and a ladder of floors by how much it keeps. P7 restates the known gap between characteristic $p$ and characteristic $0$ (the ``field with one element'' programme) as two columns of one ladder.
\item[\lab{R2}] (Old names.) P1: unique factorisation; P2: half-factorial domains (class number $2$, Carlitz), Jordan--H\"older; P3: elasticity; P5: the word problem.
\item[\lab{R3}] (Machine check, \texttt{checks/b2-p2-check.py}.) Lengths of all factorisations in $\Z[i]$ and $\Z[\sqrt{-5}]$ to size $600$ (always equal); P3.
\end{lines}

\newpage

%% ==================================================================
%%  Chapter 6  (the missing move)
%% ==================================================================
\section*{Chapter 6. The missing move}
\phantomsection\addcontentsline{toc}{section}{\texorpdfstring{Chapter 6. The missing move \quad the primes $\to$ a Survives line with no floor}{Chapter 6. The missing move}}
\begin{center}
{\normalsize the primes $\to$ a Survives line with no floor \quad$\cdot$\quad from: Chapters 3--5; Building 0, Chapters 17, 18, 34 \quad$\cdot$\quad survey}\\[4pt]
{\small The Riemann hypothesis written as a line of the grammar: a Survives line for which no floor has been built. Each attempt in this book moves the missing piece somewhere else, and every road ends at one place: an exact cancellation.}
\end{center}

\subsection*{The line}

\begin{lines}
\item[\lab{L}] (Wanted.) A floor $X$, built from $\Z$, with a move $\phi$ respecting both $+$ and $\times$, whose power ladder has a weak bound (Chapter 4, L), and an arrow back to the primes under which the eigenvalues of $\phi$ are the zeros of $\zeta$. Then Amplify gives the line. Nobody has $X$.
\end{lines}

\subsection*{Propositions: what was found, with labels}

\begin{lines}
\item[\lab{P1}] (Where the line sits, reported.) The line is equivalent to: Weil's quadratic form in the primes is $\ge 0$ on every test map (classical). Cut at size $\lambda$ (only $n \le \lambda^2$ used), the form's lowest test map has its zeros within $10^{-3}$ of the first eight zeros of $\zeta$ already at $\lambda^2 = 10$ (checked; known: Connes--Consani).
\item[\lab{P2}] (Where a bad zero hides; checked, novelty unchecked.) Plant one zero off the line at height $T$: the smallest cut at which the form turns negative is $\lambda^2 \approx T/2\pi$, almost independent of the distance from the line. A bad zero hides by height, not by nearness.
\item[\lab{P3}] (Every prime is needed; checked.) Remove one prime $p$ well inside the cut and the form drops to $-\log p / \sqrt p$; keep only finitely many primes and no size appears at all. The size is global (matches Chapter 3, T2).
\item[\lab{P4}] (A finite floor that carries the line; checked to $n = 300$, probably known.) The characteristic polynomial of the matrix $R$ of Chapter 3, P5 is $(x - 1)^{n - K - 1} Q_n(x - 1)$ with $K = \lfloor \log_2 n \rfloor$ and \[ Q_n(y) = y^{K+1} - \sum_{k=1}^{K} F_k(n)\, y^{K - k}, \] $F_k(n)$ the number of chains $1 < a_1 \mid a_2 \mid \dots \mid a_k \le n$ with proper steps: glue counts of $\times$. The line is equivalent to: the small roots have geometric mean coming near $1/\sqrt2$ at most. Random sequences with the same shape fail: it is a cancellation, not a shape.
\item[\lab{P5}] (The mixed condition; checked.) A series with a mirror but no prime product (Davenport--Heilbronn) fails the cut form at $\lambda^2 \approx 40$, caught by its zero off the line; turning its extra weights down does not repair it. Positivity needs the prime data and the mirror to belong to one object.
\end{lines}

\subsection*{Theorems}

\begin{lines}
\item[\lab{T1}] (Every road ends at cancellation.) In Building 0, Chapter 17 the orders-of-magnitude blur kills cancellation (its P6). In Chapter 4 the empty slot is a bound that loses only a constant, which needs cancellation. In P4 the condition is an alternating sum of glue counts. In P2, P3 the size is exact and global. So the missing piece is one thing seen four ways: an exact cancellation among all primes at once.
\end{lines}

\subsection*{Remarks}

\begin{lines}
\item[\lab{R1}] (Honest state.) Nothing here proves anything new about $\zeta$. The books sort the known proofs into one chain (Chapter 4), name the empty slot, and give three checked numerical facts (P2, P3, P4) whose novelty has not been checked against the literature.
\item[\lab{R2}] (Corrections kept on record.) ``$\Z$ is an island'' (no nearby prime data passes) was too strong: the passing set is convex and other weights pass at every finite cut. ``The primes are the top point of the passing set'' holds only for the first few primes ($2, 3, 5, 7, 11, 13, 19$ at $\lambda^2 \le 100$).
\item[\lab{R3}] (Old names.) L: a Hilbert--P\'olya operator, the cohomology of $\operatorname{Spec} \Z$; P1: Weil's explicit formula and positivity criterion; P5: the Davenport--Heilbronn function, the Selberg class.
\item[\lab{R4}] (Sources and checks.) The files \texttt{LOOP-RULE.md} and \texttt{amplify/} in the folder \texttt{riemann} of the project (scripts and outputs dated 11--12 Nov). P4's formula is checked in \texttt{checks/b2-p2-check.py} for $n = 6, 12, 30, 48$.
\end{lines}
